\chapter{Cơ sở lý thuyết và framework nghiên cứu}
\label{chap:theoretical-framework}

% ============================================================
% PHẦN I: LÝ THUYẾT VÀ FRAMEWORK NGHIÊN CỨU GỐC
% ============================================================

\section{Cơ sở lý thuyết và framework nghiên cứu gốc}
\label{sec:original-framework}

Phần này trình bày cơ sở phương pháp luận và framework được đề xuất
trong nghiên cứu của Eskandari et al. [1]. Trọng tâm là các nhóm
phương pháp lựa chọn đặc trưng, cơ chế kết hợp nhiều bộ chọn,
thuật toán SIFE và những thành phần hỗ trợ dự báo, diễn giải mô hình.
Các kết quả thực nghiệm do tác giả công bố được trình bày riêng
trong Chương 3.

\subsection{Dự báo chuỗi thời gian trong tiêu thụ năng lượng}
\label{subsec:time-series-theory}

Dự báo chuỗi thời gian sử dụng các quan sát trong quá khứ để ước lượng
giá trị của một hoặc nhiều thời điểm trong tương lai. Đối với bài toán
tiêu thụ năng lượng, dữ liệu đầu vào có thể bao gồm các biến đo lường
từ cảm biến trong nhà, điều kiện khí tượng, trạng thái chiếu sáng
và các thông tin liên quan đến thời gian.

Với dữ liệu đa biến, tại thời điểm $t$, hệ thống có thể được biểu diễn
bằng vector đặc trưng:

\begin{equation}
    \mathbf{x}(t)
    =
    [x_1(t),x_2(t),\ldots,x_d(t)]
    \label{eq:time-series-vector}
\end{equation}

Trong đó, $d$ là số lượng biến đầu vào. Mục tiêu dự báo là ước lượng
giá trị của biến mục tiêu tại thời điểm tương lai dựa trên những
quan sát trước đó.

Tùy theo số lượng mốc tương lai cần dự báo, bài toán có thể được chia
thành dự báo một bước (one-step-ahead forecasting) và dự báo nhiều
bước (multi-step hoặc multi-horizon forecasting). Trong dự báo một
bước, mô hình chỉ ước lượng giá trị tại thời điểm kế tiếp; trong khi
đó, dự báo nhiều bước yêu cầu mô hình tạo ra một chuỗi giá trị tương
lai tại nhiều khoảng thời gian khác nhau.

\subsection{Feature Selection và các nhóm phương pháp}
\label{subsec:feature-selection-theory}

Feature Selection là quá trình lựa chọn một tập con từ các biến đầu
vào ban đầu mà không chuyển đổi chúng sang một không gian đặc trưng
mới. Mục tiêu của quá trình này là giữ lại những biến có ích cho dự
báo, giảm số chiều dữ liệu, hạn chế thông tin dư thừa và hỗ trợ khả
năng diễn giải của mô hình.

Các phương pháp Feature Selection thường được chia thành ba nhóm
chính: Filter, Embedded và Wrapper. Ngoài ra, Hybrid Feature Selection
kết hợp ưu điểm của nhiều nhóm phương pháp, thường bằng cách sử dụng
Filter để thu hẹp tập ứng viên trước khi áp dụng Wrapper.

\small
\begin{longtable}{
    @{}
    >{\raggedright\arraybackslash}p{0.15\textwidth}
    >{\raggedright\arraybackslash}p{0.28\textwidth}
    >{\raggedright\arraybackslash}p{0.25\textwidth}
    >{\raggedright\arraybackslash}p{0.24\textwidth}
    @{}
}
\caption{So sánh các nhóm phương pháp Feature Selection.}
\label{tab:feature-selection-groups}\\

\toprule
\textbf{Nhóm}
&
\textbf{Cách đánh giá}
&
\textbf{Ưu điểm}
&
\textbf{Hạn chế chính}
\\
\midrule
\endfirsthead

\toprule
\textbf{Nhóm}
&
\textbf{Cách đánh giá}
&
\textbf{Ưu điểm}
&
\textbf{Hạn chế chính}
\\
\midrule
\endhead

\bottomrule
\endlastfoot

Filter
&
Sử dụng các thước đo thống kê hoặc thông tin giữa biến đầu vào
và biến mục tiêu.
&
Tốc độ nhanh, có thể xếp hạng nhiều biến và ít phụ thuộc vào
một estimator cụ thể.
&
Có thể bỏ qua tương tác giữa các biến hoặc không tối ưu trực tiếp
cho mô hình dự báo cuối cùng.
\\

Embedded
&
Quá trình lựa chọn đặc trưng được tích hợp trong cơ chế học
của estimator.
&
Importance phản ánh cấu trúc mà mô hình đã học và không cần
thử toàn bộ các tập con như Wrapper.
&
Kết quả phụ thuộc vào họ mô hình và cấu hình huấn luyện.
\\

Wrapper
&
Thử nghiệm các tập con đặc trưng bằng cách huấn luyện estimator
và đánh giá chất lượng dự báo.
&
Có tính model-aware, trực tiếp gắn với mục tiêu của mô hình
cuối cùng.
&
Chi phí tính toán cao và có nguy cơ overfit tập validation
nếu không kiểm soát quá trình tìm kiếm.
\\

Hybrid
&
Kết hợp Filter để thu hẹp tập ứng viên và Wrapper để lựa chọn
tập đặc trưng cuối cùng.
&
Giảm không gian tìm kiếm nhưng vẫn giữ được khả năng đánh giá
đặc trưng theo estimator.
&
Chất lượng phụ thuộc vào bước Filter; biến hữu ích có thể bị
loại trước khi Wrapper đánh giá.
\\

\end{longtable}
\normalsize

\subsection{Các phương pháp lựa chọn đặc trưng cơ sở}
\label{subsec:base-selectors-theory}

Framework nghiên cứu gốc kết hợp bốn phương pháp lựa chọn đặc trưng:
mRMR, Relief, F-test và Random Forest. Các phương pháp này sử dụng
những tiêu chí đánh giá khác nhau, qua đó cung cấp nhiều góc nhìn
về mức độ quan trọng của từng biến.

\subsubsection{Maximum Relevance Minimum Redundancy (mRMR)}

mRMR (Maximum Relevance Minimum Redundancy) sử dụng mutual information
để cân bằng hai tiêu chí: mức độ liên quan giữa một biến với biến mục
tiêu và mức độ dư thừa giữa biến đó với các biến đã được lựa chọn.

Mục tiêu là ưu tiên những biến cung cấp nhiều thông tin về target,
đồng thời hạn chế lựa chọn các biến mang thông tin gần trùng lặp
với tập đặc trưng hiện có [13].

Với $S$ là tập các biến đã được lựa chọn, điểm đánh giá của một biến
ứng viên $x_j$ có thể được biểu diễn như sau:

\begin{equation}
    \operatorname{score}(x_j)
    =
    I(x_j;y)
    -
    \frac{1}{|S|}
    \sum_{x_i \in S} I(x_j;x_i)
    \label{eq:mrmr-original}
\end{equation}

Trong đó, $I(x_j;y)$ biểu thị mutual information giữa biến $x_j$
và biến mục tiêu $y$; thành phần thứ hai đo mức độ dư thừa trung bình
giữa $x_j$ và các biến đã nằm trong tập $S$. Ở bước lựa chọn đầu tiên,
khi $S$ chưa có phần tử, thuật toán ưu tiên biến có relevance lớn nhất.

\subsubsection{Relief}

Relief là phương pháp đánh giá đặc trưng dựa trên sự khác biệt giữa
một mẫu và các mẫu lân cận gần nhất. Ý tưởng chính là một đặc trưng
có giá trị thay đổi đáng kể giữa các mẫu thuộc những nhóm khác nhau
có thể cung cấp thông tin hữu ích cho việc phân biệt các mẫu.

Phương pháp Relief cập nhật trọng số đặc trưng dựa trên mức độ khác
biệt giữa mẫu đang xét với các láng giềng gần. Các đặc trưng có khả
năng phân biệt những mẫu có target khác nhau sẽ nhận được mức đánh
giá cao hơn.

\subsubsection{F-test Feature Selection}

F-test đánh giá từng biến đầu vào thông qua thống kê $F$ của mô hình
hồi quy đơn biến. Phương pháp này đo lường mức độ liên hệ tuyến tính
giữa từng đặc trưng và biến mục tiêu.

Ưu điểm của F-test là chi phí tính toán thấp, dễ triển khai và
có khả năng cung cấp một tiêu chí thống kê độc lập với các phương
pháp còn lại. Tuy nhiên, phương pháp đánh giá từng biến riêng lẻ
nên không trực tiếp phản ánh mức độ dư thừa giữa các predictors
hoặc những quan hệ phi tuyến phức tạp.

\subsubsection{Random Forest Feature Importance}

Random Forest là mô hình ensemble gồm nhiều decision trees được
huấn luyện trên những mẫu dữ liệu và tập đặc trưng khác nhau.
Thông qua cơ chế tổng hợp nhiều cây, mô hình có khả năng biểu diễn
các quan hệ phi tuyến và tương tác giữa các biến.

Trong Feature Selection, mức độ quan trọng của một biến có thể
được đánh giá thông qua feature importance của Random Forest.
Một cách tính phổ biến là impurity-based importance, dựa trên mức
độ đóng góp của biến vào việc giảm impurity tại các nút chia.

Tuy nhiên, importance của Random Forest có thể bị ảnh hưởng bởi
cấu trúc dữ liệu, số lượng giá trị có thể chia và cấu hình mô hình.
Do đó, việc kết hợp Random Forest với những bộ chọn khác giúp
tránh phụ thuộc hoàn toàn vào một tiêu chí đánh giá duy nhất.

\subsection{Ensemble Filter}
\label{subsec:ensemble-filter-original}

Một hạn chế của việc sử dụng một bộ chọn đặc trưng đơn lẻ là kết quả
có thể phụ thuộc mạnh vào tiêu chí đánh giá của phương pháp đó.
Ví dụ, mRMR tập trung vào relevance và redundancy, F-test đánh giá
liên hệ tuyến tính từng biến, trong khi Random Forest dựa trên cấu
trúc cây quyết định.

Để tận dụng những góc nhìn bổ sung này, Eskandari et al. đề xuất
kết hợp bốn phương pháp cơ sở gồm mRMR, Relief, F-test và Random
Forest trong một Ensemble Filter [1].

Về nguyên tắc, các bộ chọn thực hiện đánh giá và xếp hạng đặc trưng
theo tiêu chí riêng. Kết quả sau đó được tổng hợp thông qua điểm
importance kết hợp và cơ chế voting. Cơ chế voting giúp đánh giá
mức độ đồng thuận giữa các bộ chọn, trong khi điểm importance
phản ánh mức độ ưu tiên của từng biến.

Framework gốc sử dụng cơ chế tổng hợp điểm và ngưỡng số phiếu
để lựa chọn những đặc trưng được nhiều phương pháp ủng hộ.
Nhờ đó, tập đặc trưng đầu ra không chỉ phụ thuộc vào thứ hạng
của một bộ chọn đơn lẻ mà còn xét đến sự đồng thuận giữa nhiều
phương pháp khác nhau.

\subsection{SIFE Wrapper và hàm mục tiêu}
\label{subsec:sife-theory}

SIFE (Set-based Integer-coded Fuzzy granular Evolutionary) là
thuật toán lựa chọn đặc trưng dựa trên biểu diễn tập hợp và tìm
kiếm tiến hóa. Phương pháp được thiết kế để tìm kiếm các tập con
đặc trưng có khả năng cân bằng giữa chất lượng dự báo và số lượng
biến được giữ lại.

Theo mô tả của nghiên cứu gốc, SIFE sử dụng biểu diễn tập hợp,
fuzzy granulation và các toán tử tiến hóa để tạo ra, đánh giá
và cập nhật quần thể nghiệm qua nhiều thế hệ. Cơ chế three-parent
crossover khai thác quan hệ union và intersection giữa các tập hợp
đặc trưng; đồng thời, fuzzy granulation hỗ trợ quá trình khởi tạo,
lựa chọn cá thể và duy trì diversity trong quần thể [1,4].

Hàm mục tiêu của thuật toán kết hợp prediction error và tỷ lệ
số lượng đặc trưng được lựa chọn:

\begin{equation}
    F(S)
    =
    w_1 \cdot \operatorname{Er}(S)
    +
    w_2 \cdot \operatorname{Ps}(S)
    \label{eq:sife-objective-original}
\end{equation}

Trong đó:

\begin{itemize}
    \item $\operatorname{Er}(S)$ là prediction error của estimator
    khi sử dụng tập đặc trưng $S$.
    \item $\operatorname{Ps}(S)$ là tỷ lệ số đặc trưng được lựa chọn
    trên tổng số đặc trưng ứng viên.
    \item $w_1$ và $w_2$ là các trọng số điều chỉnh mức độ ưu tiên
    giữa chất lượng dự báo và tính gọn nhẹ của tập đặc trưng.
\end{itemize}

Nghiên cứu gốc sử dụng $w_1=0{,}80$ và $w_2=0{,}20$. Như vậy,
thuật toán ưu tiên chất lượng dự báo nhưng vẫn xem xét mức độ
gọn nhẹ của tập đặc trưng. Một tập con có ít biến hơn có thể
giảm thành phần phạt về kích thước, nhưng nếu loại bỏ những
biến quan trọng thì prediction error có thể tăng.

Các tham số SIFE được mô tả trong bài báo gồm population size
bằng 120, số vòng lặp tối đa bằng 60, crossover probability
bằng 0,6 và mutation probability bằng 0,5 [1].

\subsection{Các chiến lược Feature Selection trong framework gốc}
\label{subsec:selection-strategies-original}

Framework nghiên cứu kết hợp các phương pháp lựa chọn đặc trưng
thành ba chiến lược chính, đồng thời sử dụng trường hợp không lựa
chọn đặc trưng làm baseline tham chiếu.

\begin{itemize}
    \item \textbf{S0 -- No Feature Selection:} sử dụng toàn bộ
    đặc trưng đầu vào, không thực hiện bước lựa chọn biến.
    Đây là baseline để đánh giá ảnh hưởng của Feature Selection.

    \item \textbf{S1 -- Ensemble Filter:} sử dụng Ensemble Filter
    để lựa chọn tập đặc trưng dựa trên kết quả tổng hợp của các
    bộ chọn cơ sở.

    \item \textbf{S2 -- SIFE Wrapper:} sử dụng trực tiếp thuật toán
    SIFE trên tập đặc trưng ban đầu. Quá trình tìm kiếm được đánh
    giá theo estimator cụ thể.

    \item \textbf{S1S2 -- Hybrid Filter-Wrapper:} kết hợp Ensemble
    Filter và SIFE. Ensemble Filter thu hẹp tập đặc trưng ứng viên,
    sau đó SIFE tiếp tục tìm kiếm tập con phù hợp với estimator.
\end{itemize}

Sự khác biệt về đầu vào của các chiến lược được thể hiện trong
Bảng~\ref{tab:original-selection-strategies}.

\small
\begin{longtable}{
    @{}
    >{\raggedright\arraybackslash}p{0.13\textwidth}
    >{\raggedright\arraybackslash}p{0.35\textwidth}
    >{\raggedright\arraybackslash}p{0.39\textwidth}
    @{}
}
\caption{Cấu trúc các chiến lược Feature Selection trong framework gốc.}
\label{tab:original-selection-strategies}\\

\toprule
\textbf{Chiến lược}
&
\textbf{Quy trình lựa chọn}
&
\textbf{Mục đích}
\\
\midrule
\endfirsthead

\toprule
\textbf{Chiến lược}
&
\textbf{Quy trình lựa chọn}
&
\textbf{Mục đích}
\\
\midrule
\endhead

\bottomrule
\endlastfoot

S0
&
Sử dụng toàn bộ đặc trưng ban đầu.
&
Tạo baseline để so sánh với các chiến lược có lựa chọn đặc trưng.
\\

S1
&
Áp dụng Ensemble Filter trên tập đặc trưng ban đầu.
&
Lựa chọn đặc trưng dựa trên sự kết hợp của nhiều bộ chọn cơ sở.
\\

S2
&
Áp dụng SIFE trực tiếp trên tập đặc trưng ban đầu.
&
Tìm kiếm tập con theo chất lượng dự báo của estimator.
\\

S1S2
&
Áp dụng Ensemble Filter trước, sau đó sử dụng SIFE trên tập
đặc trưng đã được rút gọn.
&
Kết hợp khả năng lọc nhanh của Filter và khả năng tối ưu
theo estimator của Wrapper.
\\

\end{longtable}
\normalsize

\subsection{Các mô hình dự báo}
\label{subsec:forecasting-models-original}

Sau bước lựa chọn đặc trưng, framework sử dụng các mô hình dự báo
để đánh giá hiệu quả của từng tập đặc trưng. Nghiên cứu gốc xem xét
nhiều họ mô hình khác nhau, từ hồi quy tuyến tính đến các mô hình
phi tuyến và mô hình dựa trên kernel.

\subsubsection{LASSO}

LASSO (Least Absolute Shrinkage and Selection Operator) là mô hình
hồi quy tuyến tính sử dụng regularization L1. Thành phần phạt L1
giúp co một số hệ số hồi quy về 0, qua đó tạo ra mô hình thưa và
có thể hỗ trợ lựa chọn biến.

\subsubsection{SVR-Linear}

Support Vector Regression với linear kernel xây dựng hàm hồi quy
trong không gian đặc trưng tuyến tính. Mô hình phù hợp khi quan hệ
giữa đầu vào và target có thể được biểu diễn tương đối tốt bằng
một hàm tuyến tính.

\subsubsection{SVR-RBF}

SVR với Radial Basis Function kernel có khả năng biểu diễn các
quan hệ phi tuyến thông qua ánh xạ kernel. Hiệu quả của mô hình
phụ thuộc vào các tham số như $C$, $\gamma$ và $\epsilon$.

\subsubsection{Gaussian Process Regression}

Gaussian Process Regression (GPR) là phương pháp hồi quy xác suất
dựa trên Gaussian Process. Mô hình sử dụng kernel để biểu diễn quan
hệ giữa các quan sát và có thể cung cấp phân phối dự báo, qua đó
hỗ trợ đánh giá mức độ không chắc chắn.

\subsubsection{LSBoost}

LSBoost là phương pháp gradient boosting sử dụng squared error
làm hàm mất mát. Các weak learners, thường là decision trees,
được xây dựng tuần tự để học phần sai số còn lại của tổ hợp mô hình
trước đó. Cơ chế này cho phép mô hình biểu diễn các quan hệ phi
tuyến trong dữ liệu.

\subsection{Bayesian Optimization}
\label{subsec:bayesian-optimization-original}

Bayesian Optimization là phương pháp tối ưu hóa các hàm mục tiêu
có chi phí đánh giá cao. Thay vì thử toàn bộ không gian tham số,
phương pháp sử dụng surrogate model để xấp xỉ hàm mục tiêu và
acquisition function để đề xuất cấu hình tiếp theo cần đánh giá.

Trong framework nghiên cứu, Bayesian Optimization được sử dụng
để tìm kiếm cấu hình tham số phù hợp cho các mô hình dự báo.
Cách tiếp cận này giúp khai thác thông tin từ những lần đánh giá
trước đó và giảm số lần thử nghiệm so với việc tìm kiếm vét cạn
trên toàn bộ không gian tham số [1,10].

\subsection{SHAP và khả năng diễn giải mô hình}
\label{subsec:shap-theory}

SHAP (SHapley Additive exPlanations) dựa trên giá trị Shapley
trong cooperative game theory để phân bổ phần đóng góp của từng
đặc trưng vào dự báo của mô hình.

Đối với một mẫu đầu vào $x$, biểu diễn giải thích có dạng:

\begin{equation}
    f(x) = \phi_0 + \sum_{j=1}^{M}\phi_j
    \label{eq:shap-additive-original}
\end{equation}

Trong đó, $\phi_0$ là giá trị cơ sở của mô hình và $\phi_j$
là phần đóng góp của đặc trưng thứ $j$ vào dự báo tại mẫu đang xét.

Giá trị tuyệt đối trung bình của SHAP trên nhiều mẫu thường được
sử dụng để đánh giá global feature importance. Ngoài ra, SHAP
có thể được trình bày dưới dạng summary plot, dependence plot
hoặc heatmap để quan sát mức độ và chiều tác động của các biến.

Trong nghiên cứu gốc, SHAP được sử dụng như một lớp phân tích
bổ sung nhằm hỗ trợ diễn giải các mô hình dự báo và vai trò của
những đặc trưng được lựa chọn [1,8].

% ============================================================
% PHẦN II: FRAMEWORK ĐƯỢC HIỆN THỰC HÓA TRONG NGHIÊN CỨU
% ============================================================

\section{Framework được hiện thực hóa trong nghiên cứu này}
\label{sec:implemented-framework}

Dựa trên framework nghiên cứu gốc, nghiên cứu này xây dựng một
phiên bản triển khai bằng Python trên bộ dữ liệu KAG Energy
Consumption. Mục tiêu là giữ lại logic so sánh giữa các chiến
lược lựa chọn đặc trưng, đồng thời điều chỉnh một số thuật toán
và công cụ để phù hợp với dữ liệu, thư viện và môi trường thực
nghiệm.

Do khác biệt về bộ dữ liệu và môi trường triển khai, phiên bản
này không được xem là bản sao hoàn toàn của mã nguồn MATLAB
trong nghiên cứu gốc. Những điểm giữ nguyên, điều chỉnh và
thay thế được trình bày cụ thể trong các mục dưới đây.

\subsection{Định nghĩa bài toán dự báo}
\label{subsec:implemented-forecasting-task}

Nghiên cứu sử dụng dữ liệu tiêu thụ năng lượng trong bộ dữ liệu
KAG Energy Consumption. Dữ liệu đầu vào có dạng đa biến, bao gồm
các biến liên quan đến mức tiêu thụ năng lượng, cảm biến môi
trường trong nhà, điều kiện khí tượng, trạng thái chiếu sáng
và thông tin thời gian.

Khác với bài toán dự báo nhiều bước, notebook xây dựng bài toán
hồi quy dự báo một bước. Biến mục tiêu là giá trị tiêu thụ năng
lượng của thiết bị tại thời điểm kế tiếp:

\begin{equation}
    y(t) = \texttt{Appliances}(t+1)
    \label{eq:implemented-target}
\end{equation}

Các biến đầu vào được sử dụng để ước lượng giá trị target tại
thời điểm $t+1$. Do đó, đây là bài toán one-step-ahead regression,
không phải bài toán dự báo đồng thời nhiều horizon.

\subsection{Hiện thực hóa các base selectors}
\label{subsec:implemented-base-selectors}

Nghiên cứu sử dụng bốn base selectors tương ứng với cấu trúc
Ensemble Filter trong framework gốc. Tuy nhiên, một số phương pháp
được hiện thực hóa thông qua thư viện Python phù hợp với đặc điểm
của target liên tục.

\subsubsection{mRMR bằng Mutual Information Regression}

Notebook sử dụng \texttt{mutual\_info\_regression} để ước lượng
mutual information giữa từng biến đầu vào và target liên tục.
Quá trình lựa chọn được triển khai theo thủ tục greedy:
ở bước đầu, thuật toán ưu tiên biến có relevance cao nhất;
ở các bước tiếp theo, điểm relevance được điều chỉnh bằng
mức redundancy trung bình với những biến đã lựa chọn.

Đây là cách triển khai mRMR bằng Python cho bài toán regression,
không phải sao chép trực tiếp mã nguồn MATLAB của nghiên cứu gốc.

\subsubsection{Random Forest Importance}

Random Forest được triển khai bằng
\texttt{RandomForestRegressor}. Notebook sử dụng cấu hình gồm
250 cây, \texttt{max\_features=0.8} và
\texttt{min\_samples\_leaf=2}.

Mức độ quan trọng của biến được lấy từ thuộc tính
\texttt{feature\_importances\_}. Đây là impurity-based importance,
do đó kết quả có thể chịu ảnh hưởng bởi cấu trúc biến và cơ chế
chia nhánh của decision trees.

\subsubsection{RReliefF cho target liên tục}

Relief trong framework gốc được điều chỉnh trong notebook thành
RReliefF để phù hợp với bài toán hồi quy có target liên tục [5].

RReliefF đánh giá đặc trưng thông qua quan hệ giữa sự thay đổi
của giá trị đặc trưng và sự thay đổi của target giữa những mẫu
lân cận. Việc sử dụng RReliefF là một adaptation phù hợp với
bản chất regression của bài toán KAG; tuy nhiên, đây không phải
cùng một implementation Relief classification được trình bày
trong nghiên cứu gốc.

\subsubsection{F-test Feature Selection}

F-test được sử dụng để xếp hạng các biến theo thống kê $F$
của hồi quy đơn biến. Phương pháp cung cấp một góc nhìn tuyến tính
và có chi phí tính toán thấp, bổ sung cho các bộ chọn dựa trên
mutual information, láng giềng và cây quyết định.

\subsection{Hiện thực hóa Ensemble Filter}
\label{subsec:implemented-ensemble-filter}

Trong notebook, mỗi base selector thực hiện xếp hạng và đề xuất
tối đa $m=10$ đặc trưng từ 36 candidate features. Kết quả của
bốn bộ chọn được tổng hợp dựa trên hai thành phần: điểm importance
và số lượng phiếu ủng hộ.

Cụ thể, \texttt{vote\_count} đếm số base selectors đưa một biến
vào danh sách Top 10. Các giá trị importance được chuẩn hóa về
khoảng $[0,1]$; riêng mRMR sử dụng relative rank để đưa điểm
xếp hạng về cùng thang đo. Sau đó, điểm trung bình được tính
với trọng số đồng đều $0{,}25$ cho mỗi bộ chọn.

Công thức tổng hợp được biểu diễn như sau:

\begin{equation}
    \operatorname{AvgImp}(x)
    =
    \frac{1}{4}\sum_{k=1}^{4} I_k(x)
    \label{eq:implemented-avg-importance}
\end{equation}

\begin{equation}
    \operatorname{Vote}(x)
    =
    \sum_{k=1}^{4}
    \mathbb{1}
    \left\{
        x \in \operatorname{Top10}(k)
    \right\}
    \label{eq:implemented-voting}
\end{equation}

Trong đó, $I_k(x)$ là importance đã chuẩn hóa của biến $x$
theo bộ chọn thứ $k$; $\mathbb{1}\{\cdot\}$ là hàm chỉ báo,
nhận giá trị 1 nếu điều kiện đúng và 0 nếu điều kiện sai.

Điều kiện lựa chọn chính là:

\begin{equation}
    \operatorname{Vote}(x) \geq 2
    \label{eq:implemented-vote-threshold}
\end{equation}

Những biến thỏa mãn điều kiện trên được sắp xếp theo
$\operatorname{AvgImp}$ và lựa chọn tối đa 10 biến có điểm
cao nhất.

Cách triển khai này giúp việc tổng hợp importance và voting
trở nên minh bạch, dễ kiểm tra. Tuy nhiên, nó không hoàn toàn
giống cơ chế elimination và weighted voting được mô tả trong
nghiên cứu gốc. Vì vậy, kết quả S1 trong nghiên cứu này cần
được hiểu là kết quả của phiên bản Ensemble Filter đã được
hiện thực hóa trong notebook.

\subsection{Hiện thực hóa SIFE bằng Python}
\label{subsec:implemented-sife}

Notebook xây dựng một Python port của SIFE dựa trên mô tả
thuật toán trong nghiên cứu gốc. Phiên bản này giữ lại hàm
mục tiêu và biểu diễn tập hợp đặc trưng, đồng thời triển khai
các toán tử tìm kiếm tiến hóa bằng những cấu trúc phù hợp
với môi trường Python.

Hàm mục tiêu được sử dụng là:

\begin{equation}
    F(S)
    =
    0{,}80 \cdot \operatorname{Er}(S)
    +
    0{,}20 \cdot \operatorname{Ps}(S)
    \label{eq:implemented-sife-objective}
\end{equation}

Trong đó, $\operatorname{Er}(S)$ là prediction error trên
validation và $\operatorname{Ps}(S)$ là tỷ lệ số đặc trưng
được chọn trên tổng số đặc trưng ứng viên.

Các thành phần chính trong Python port gồm:

\begin{itemize}
    \item \textbf{Biểu diễn nghiệm:} mỗi cá thể được biểu diễn
    bằng một tập hợp các feature được lựa chọn.

    \item \textbf{Khởi tạo quần thể:} các nghiệm ban đầu được
    khởi tạo có phân tầng theo cardinality để tạo ra những
    tập con có kích thước khác nhau.

    \item \textbf{Duy trì diversity:} khoảng cách Jaccard giữa
    các tập hợp được sử dụng như một proxy để đánh giá mức độ
    khác biệt giữa các nghiệm. Đây là phương án thay thế cho
    fuzzy granulation nguyên bản.

    \item \textbf{Lựa chọn cá thể:} thuật toán sử dụng cơ chế
    lựa chọn elite kết hợp tournament selection.

    \item \textbf{Crossover:} tournament selection lựa chọn
    ba parent; các phép union và intersection được sử dụng
    để tạo ra nghiệm con.

    \item \textbf{Mutation:} phép bù tập hợp được sử dụng để
    thêm hoặc loại bỏ đặc trưng khỏi nghiệm hiện tại.
\end{itemize}

Python port giữ các tham số population size bằng 120,
maximum iteration bằng 60, crossover probability bằng 0,6
và mutation probability bằng 0,5 theo cấu hình được mô tả
trong nghiên cứu gốc.

Tuy nhiên, cần lưu ý rằng fuzzy granulation nguyên bản không
được triển khai. Khoảng cách Jaccard chỉ đóng vai trò proxy
diversity, do đó cơ chế tìm kiếm không hoàn toàn tương đương
với SIFE gốc. Ngoài ra, notebook sử dụng vòng lặp
\texttt{range(max\_iter + 1)}, tạo các mốc generation từ 0
đến 60, tức 61 mốc được ghi nhận. Đây là một khác biệt nhỏ
về cách đánh số vòng lặp so với cách mô tả 60 iterations
trong bài báo.

Vì vậy, SIFE trong nghiên cứu này được xem là một phiên bản
hiện thực hóa dựa trên mô tả thuật toán, không phải bản sao
nguyên bản của source code MATLAB.

\subsection{Thiết kế các nhánh thực nghiệm S0, S1, S2 và S1S2}
\label{subsec:implemented-selection-strategies}

Dựa trên cấu trúc framework gốc, nghiên cứu xây dựng bốn nhánh
thực nghiệm để so sánh ảnh hưởng của các chiến lược lựa chọn
đặc trưng trên cùng bài toán dự báo.

\begin{itemize}
    \item \textbf{S0 -- Baseline:} sử dụng toàn bộ 36 candidate
    features mà không thực hiện Feature Selection.

    \item \textbf{S1 -- Ensemble Filter:} áp dụng Ensemble Filter
    trên toàn bộ 36 candidate features. Tập đặc trưng được chọn
    sau đó được sử dụng làm đầu vào cho các mô hình dự báo.

    \item \textbf{S2 -- Wrapper:} áp dụng SIFE trực tiếp trên
    toàn bộ 36 candidate features. Quá trình tìm kiếm được
    thực hiện riêng theo từng estimator.

    \item \textbf{S1S2 -- Hybrid Filter-Wrapper:} áp dụng SIFE
    độc lập trên tập đặc trưng đã được rút gọn bởi S1. Mỗi
    estimator thực hiện quá trình tìm kiếm trên cùng không
    gian ứng viên S1 nhưng có thể đưa ra tập con khác nhau.
\end{itemize}

Điểm quan trọng trong thiết kế là S1S2 không được tạo bằng cách
lấy giao hậu nghiệm giữa kết quả S1 và S2. Thay vào đó, SIFE
được khởi chạy độc lập trên tập ứng viên S1. Đây là cách hiện
thực hóa trực tiếp logic Hybrid Filter-Wrapper: Filter thực
hiện bước thu hẹp không gian, sau đó Wrapper tiếp tục lựa
chọn đặc trưng theo estimator.

Bảng~\ref{tab:implemented-selection-strategies} tóm tắt
đầu vào và vai trò của từng nhánh trong thực nghiệm.

\small
\begin{longtable}{
    @{}
    >{\raggedright\arraybackslash}p{0.12\textwidth}
    >{\raggedright\arraybackslash}p{0.29\textwidth}
    >{\raggedright\arraybackslash}p{0.26\textwidth}
    >{\raggedright\arraybackslash}p{0.25\textwidth}
    @{}
}
\caption{Thiết kế các nhánh Feature Selection trong nghiên cứu này.}
\label{tab:implemented-selection-strategies}\\

\toprule
\textbf{Nhánh}
&
\textbf{Đầu vào}
&
\textbf{Phương pháp}
&
\textbf{Vai trò thực nghiệm}
\\
\midrule
\endfirsthead

\toprule
\textbf{Nhánh}
&
\textbf{Đầu vào}
&
\textbf{Phương pháp}
&
\textbf{Vai trò thực nghiệm}
\\
\midrule
\endhead

\bottomrule
\endlastfoot

S0
&
36 candidate features.
&
Không thực hiện Feature Selection.
&
Baseline để đánh giá hiệu quả của việc lựa chọn đặc trưng.
\\

S1
&
36 candidate features.
&
Ensemble Filter.
&
Đánh giá chiến lược Filter kết hợp nhiều bộ chọn.
\\

S2
&
36 candidate features.
&
SIFE, chạy riêng theo từng estimator.
&
Đánh giá chiến lược Wrapper trực tiếp.
\\

S1S2
&
Tập đặc trưng đầu ra từ S1.
&
SIFE, chạy riêng theo từng estimator.
&
Đánh giá chiến lược Hybrid Filter-Wrapper.
\\

\end{longtable}
\normalsize

\subsection{Các mô hình dự báo trong thực nghiệm}
\label{subsec:implemented-models}

Nghiên cứu sử dụng năm mô hình dự báo để đánh giá các nhánh
Feature Selection: LASSO, SVR-Linear, SVR-RBF, GPR và LSBoost.

Các mô hình được lựa chọn nhằm tạo ra sự đa dạng về cơ chế học,
bao gồm hồi quy tuyến tính có regularization, hồi quy dựa trên
support vector, mô hình Gaussian Process và mô hình ensemble
dựa trên boosting.

Trong notebook, các mô hình được hiện thực hóa bằng thư viện
scikit-learn:

\begin{itemize}
    \item \textbf{LASSO:} sử dụng
    \texttt{sklearn.linear\_model.Lasso}, với regularization L1
    để kiểm soát hệ số hồi quy.

    \item \textbf{SVR-Linear:} sử dụng
    \texttt{sklearn.svm.SVR} với linear kernel.

    \item \textbf{SVR-RBF:} sử dụng
    \texttt{sklearn.svm.SVR} với RBF kernel để biểu diễn
    các quan hệ phi tuyến.

    \item \textbf{GPR:} sử dụng
    \texttt{GaussianProcessRegressor} với kernel kết hợp
    \texttt{ConstantKernel}, \texttt{RBF} và
    \texttt{WhiteKernel}. Số hàng huấn luyện được giới hạn
    để kiểm soát chi phí tính toán.

    \item \textbf{LSBoost:} được ánh xạ sang
    \texttt{GradientBoostingRegressor} với loss
    \texttt{squared\_error}.
\end{itemize}

Việc sử dụng các implementation trong scikit-learn giúp thống
nhất môi trường thực nghiệm và thuận tiện cho quá trình huấn
luyện, tối ưu tham số và đánh giá mô hình.

\subsection{Hiện thực hóa Bayesian Optimization}
\label{subsec:implemented-bayesian-optimization}

Thay vì sử dụng công cụ Bayesian Optimization trong MATLAB
như nghiên cứu gốc, notebook sử dụng
\texttt{BayesSearchCV} từ thư viện scikit-optimize.

Quá trình tìm kiếm tham số được kết hợp với
\texttt{TimeSeriesSplit} gồm ba folds để duy trì thứ tự thời
gian trong quá trình đánh giá. Mỗi mô hình sử dụng 20 iterations
tìm kiếm cấu hình.

Việc tối ưu tham số chỉ được thực hiện trên training partition.
Holdout test được giữ riêng để đánh giá cuối cùng và không
tham gia vào quá trình lựa chọn tham số. Cách tổ chức này
nhằm hạn chế việc sử dụng thông tin từ tập kiểm tra trong
quá trình tối ưu mô hình.

\subsection{Hiện thực hóa SHAP}
\label{subsec:implemented-shap}

Notebook sử dụng \texttt{PermutationExplainer} để ước lượng
đóng góp của các đặc trưng vào dự báo. Khác với việc tính toán
đầy đủ mọi tổ hợp đặc trưng để thu được Shapley values chính
xác, PermutationExplainer sử dụng các phép hoán vị để xấp xỉ
giá trị đóng góp.

Cấu hình giải thích trong notebook gồm 40 background rows
và 30 rows được đưa vào quá trình giải thích. Ngân sách đánh
giá được đặt theo công thức:

\begin{equation}
    \texttt{max\_evals}=2M+1
    \label{eq:shap-max-evals}
\end{equation}

Trong đó, $M$ là số đặc trưng đầu vào của mô hình đang được
giải thích. Vì vậy, ngân sách đánh giá thay đổi theo số lượng
đặc trưng của từng nhánh.

Do sử dụng PermutationExplainer với số mẫu và ngân sách đánh
giá giới hạn, kết quả SHAP trong nghiên cứu này là kết quả
xấp xỉ. Các giá trị importance được sử dụng để hỗ trợ diễn
giải hành vi mô hình, không được xem là bằng chứng về quan
hệ nhân quả giữa biến đầu vào và mức tiêu thụ năng lượng.

\subsection{Tổng hợp những điểm điều chỉnh so với framework gốc}
\label{subsec:framework-deviations}

Mặc dù nghiên cứu giữ lại cấu trúc phương pháp luận chính
của framework gốc, một số thành phần được điều chỉnh để phù
hợp với dữ liệu KAG và môi trường Python. Các khác biệt này
cần được xác định rõ để tránh diễn giải kết quả thực nghiệm
như thể hai nghiên cứu sử dụng hoàn toàn cùng một quy trình.

\small
\begin{longtable}{
    @{}
    >{\raggedright\arraybackslash}p{0.19\textwidth}
    >{\raggedright\arraybackslash}p{0.35\textwidth}
    >{\raggedright\arraybackslash}p{0.35\textwidth}
    @{}
}
\caption{Đối chiếu framework nghiên cứu gốc và phiên bản triển khai.}
\label{tab:framework-comparison}\\

\toprule
\textbf{Thành phần}
&
\textbf{Nghiên cứu gốc}
&
\textbf{Nghiên cứu này}
\\
\midrule
\endfirsthead

\toprule
\textbf{Thành phần}
&
\textbf{Nghiên cứu gốc}
&
\textbf{Nghiên cứu này}
\\
\midrule
\endhead

\bottomrule
\endlastfoot

Dữ liệu
&
Sử dụng bộ dữ liệu tiêu thụ năng lượng trong nghiên cứu gốc.
&
Sử dụng bộ dữ liệu KAG Energy Consumption.
\\

Bài toán dự báo
&
Theo thiết kế thực nghiệm của nghiên cứu gốc.
&
Hồi quy dự báo một bước với target
\texttt{Appliances}$(t+1)$.
\\

mRMR
&
Sử dụng mRMR trong nhóm base selectors.
&
Triển khai bằng \texttt{mutual\_info\_regression}
và thủ tục greedy cho target liên tục.
\\

Relief
&
Sử dụng Relief trong Ensemble Filter.
&
Sử dụng RReliefF để phù hợp với bài toán regression.
\\

Random Forest
&
Sử dụng Random Forest trong nhóm base selectors.
&
Sử dụng \texttt{RandomForestRegressor} và
impurity-based importance.
\\

Ensemble Filter
&
Tổng hợp importance và voting theo framework gốc.
&
Chuẩn hóa importance về $[0,1]$, lấy trung bình đồng trọng số
và áp dụng ngưỡng voting tối thiểu hai phiếu.
\\

SIFE
&
Sử dụng set-based representation, fuzzy granulation
và các toán tử tiến hóa.
&
Triển khai Python port; dùng Jaccard distance làm proxy
diversity thay cho fuzzy granulation nguyên bản.
\\

Bayesian Optimization
&
Sử dụng công cụ Bayesian Optimization trong MATLAB.
&
Sử dụng \texttt{BayesSearchCV}, kết hợp
\texttt{TimeSeriesSplit} ba folds và 20 iterations.
\\

SHAP
&
Sử dụng SHAP để phân tích khả năng diễn giải của mô hình.
&
Sử dụng \texttt{PermutationExplainer} với ngân sách đánh giá
giới hạn.
\\

\end{longtable}
\normalsize

Những điều chỉnh trên có ý nghĩa quan trọng khi phân tích kết quả.
Thứ nhất, kết quả trên KAG phản ánh khả năng hoạt động của framework
được hiện thực hóa trong môi trường Python, không phải một phép
tái lập hoàn toàn giống nghiên cứu gốc. Thứ hai, các thay đổi như
RReliefF, cơ chế tổng hợp Ensemble Filter và proxy diversity của
SIFE có thể ảnh hưởng đến tập đặc trưng cuối cùng. Thứ ba, sự
khác biệt về dữ liệu và thiết kế bài toán cũng khiến các chỉ số
đánh giá giữa hai nghiên cứu không thể được so sánh trực tiếp
một cách tuyệt đối.

\subsection{Vai trò của framework tích hợp trong nghiên cứu}
\label{subsec:integrated-framework-role}

Framework tích hợp được xây dựng nhằm đánh giá đồng thời ba khía
cạnh: chất lượng dự báo, mức độ thu gọn tập đặc trưng và khả năng
diễn giải mô hình.

Ensemble Filter kết hợp nhiều tiêu chí đánh giá để giảm sự phụ
thuộc vào một bộ chọn đơn lẻ. SIFE tiếp tục đánh giá các tập con
theo prediction error của estimator, đồng thời xem xét số lượng
đặc trưng được giữ lại. Nhánh Hybrid S1S2 kết hợp hai bước này
bằng cách thu hẹp tập ứng viên trước khi thực hiện tìm kiếm
Wrapper.

Về mặt không gian tìm kiếm, nếu số lượng biến ứng viên giảm từ
36 xuống 10, số tập con tiềm năng giảm từ $2^{36}$ xuống $2^{10}$.
Đây là lợi ích về quy mô không gian tìm kiếm, không có nghĩa
thuật toán SIFE đã duyệt toàn bộ các tập con có thể có. SIFE
vẫn là một thuật toán tìm kiếm heuristic với số lượng vòng lặp
hữu hạn.

Cuối cùng, SHAP cung cấp thông tin bổ sung về mức độ đóng góp
của các biến đối với dự báo của mô hình. Nhờ đó, framework
không chỉ đánh giá mô hình dựa trên sai số dự báo mà còn hỗ trợ
phân tích vai trò của tập đặc trưng được lựa chọn.

Như vậy, framework được hiện thực hóa trong nghiên cứu này
kế thừa logic chính của nghiên cứu gốc nhưng có những điều
chỉnh cụ thể về thuật toán và công cụ. Các kết quả thực nghiệm
trên bộ dữ liệu KAG sẽ được trình bày ở những chương tiếp theo
để đánh giá hiệu quả của từng chiến lược trong điều kiện dữ
liệu và triển khai đã lựa chọn.